% !TeX root = ../thuthesis-example.tex

\chapter{绪论}

\section{国内研究背景}

随着人工智能技术在教育信息化中的应用不断加深，学习行为分析、个性化推荐和智能反馈已经成为智慧教学平台的重要组成部分。在国内高校和在线教育场景中，教师不再只关注学生是否完成任务，也越来越关注学生在学习过程中的答题轨迹、错误类型、练习频率和能力变化。对于国际中文教育而言，拼音、声调、词汇和听读训练具有明显的层级依赖关系，学生的错误往往不是孤立发生，而是与前置知识点掌握不足有关。

目前许多中文学习平台已经具备课程发布、作业提交和学习记录展示能力，但真正能把学习数据转化为可解释推荐结果的系统仍然不足。尤其在本地演示和课程设计场景中，若完全依赖大型模型或外部 API，容易受到网络、成本和可解释性限制。因此，采用 DKT 知识追踪、编辑距离校准和轻量推荐策略，能够在工程可控的前提下体现人工智能方法的实际价值。

\section{国外研究背景}

在国际研究中，Deep Knowledge Tracing（DKT）是教育数据挖掘领域的经典方法之一。该方法通过序列模型刻画学生知识状态随时间变化的过程，能够根据历史答题序列预测学生对后续知识点的掌握概率。与静态正确率统计相比，知识追踪更适合处理连续学习过程中的动态能力变化。

在语言学习和文本处理领域，Levenshtein Distance 是经典字符串相似度算法，常用于拼写纠错、语音识别后处理和文本对齐。对于中文学习中的拼音、声调和字幕转写任务，编辑距离可以以较低计算开销发现相近发音或相近拼写偏误，并为系统生成可解释的校准证据。将知识追踪和文本校准结合，可以使推荐结果既有学生掌握状态依据，也有具体语言偏误依据。

\section{研究目的}

本研究围绕 LingoBridge 国际中文学习平台，设计并实现“基于 DKT 知识追踪与 Levenshtein 距离校准的国际中文自适应学习推荐系统”。项目的目标是在不破坏现有教师端、学生端和管理端 MVP 主流程的基础上，增加一个可本地运行的人工智能分析模块，使平台具备学习状态诊断、字幕或拼音校准、薄弱知识点识别和个性化练习推荐能力。

具体而言，本项目希望完成三项任务：第一，构建学习行为数据处理流程，对课程、作业、录音练习和反馈记录进行清洗与特征构建；第二，使用 DKT 模型预测学生对拼音、声调和词汇知识点的掌握状态，并用 Levenshtein 距离校准 ASR 字幕或拼音转写偏误；第三，将算法结果反馈到学生端、教师端和管理端，使系统从教学管理工具进一步扩展为具备人工智能辅助决策能力的学习平台。
